% ============================================================
%  INTRODUCTION
% ============================================================
\section{Introduction}

The human auditory system is remarkably adept at extracting spatial information from sound —
identifying the direction, distance, and movement of sources in three-dimensional space.
This perceptual capability, known as \textit{sound spatialisation}, has become one of the most
actively researched and commercially significant frontiers in modern audio engineering.
As digital experiences grow increasingly immersive, the faithful reproduction of spatial acoustic
cues through headphones and loudspeakers is no longer a niche pursuit but a mainstream
requirement spanning entertainment, professional audio, and accessibility.

The commercial momentum behind spatial audio is substantial. The global immersive entertainment
market is projected to reach \$474 billion by 2030, rising sharply from \$133.6 billion in
2024~\cite{immersive_market_2024}. This rapid expansion is fuelling demand for spatial audio
technologies across gaming, music production, and interactive media. Major platforms have already
responded: Apple Music introduced Spatial Audio in June 2021 and has since offered increased
royalty rates to incentivise artists to deliver content in immersive formats~\cite{apple_spatial_royalties},
signalling that spatialised sound is becoming a standard expectation rather than a premium feature.

Consumer adoption mirrors industry investment. A recent survey found that nearly 80\% of
Generation Z respondents would purchase new headphones specifically for spatial audio
capabilities, with more than one quarter willing to do so regardless of cost~\cite{genz_survey}.
These figures underscore a generational shift in how audiences perceive and prioritise
three-dimensional audio experiences.

The impact is particularly pronounced in two professional communities that stand to benefit most
from precise, reliable spatialisation: \textbf{competitive gamers} and \textbf{musicians}.

For professional esports athletes and competitive gamers, spatial audio confers measurable
tactical advantages. Accurate sound localisation enables players to determine the precise
direction and distance of in-game events — footsteps, projectiles, environmental cues — before
they are visible on screen. Beyond competitive advantage, spatialised game audio enhances
atmospheric immersion, sharpens spatial awareness, and can improve accessibility for players
with visual impairments by providing richer non-visual feedback~\cite{game_audio_benefits}.
In high-stakes esports environments, where reaction times are measured in milliseconds,
the fidelity of spatial audio rendering can directly influence performance outcomes.

For musicians, audio engineers, and music producers, the ability to accurately perceive and
evaluate spatial audio is equally critical. As immersive formats such as Dolby Atmos and
Sony 360 Reality Audio become production standards, professionals must develop a reliable
perceptual reference for three-dimensional sound placement. Errors in spatial judgement
during mixing translate directly into poor listener experiences across playback systems.
There is therefore a pressing need for standardised, reproducible tools that can measure
and train an individual's spatial hearing acuity.

This project addresses that need by designing, implementing, and evaluating a web-based
platform for spatial hearing assessment. The system combines a rigorously designed
psychoacoustic test protocol with a modern, accessible web application, enabling users
to quantify their spatialisation perception across key auditory dimensions. By targeting
both gamers and audio professionals, the platform occupies a unique position at the
intersection of psychoacoustics, human–computer interaction, and immersive audio technology.

The remainder of this report is structured as follows. Section~\ref{sec:wp2} details the
synthesis and selection of the spatialisation test protocol (WP2). Section~\ref{sec:wp4}
describes the architecture and development of the web platform (WP4). Subsequent sections
present experimental results, discussion, and conclusions.

% ============================================================
%  WP2 — SYNTHESIS OF AUDIOMETRIC SPATIALISATION TEST
% ============================================================
\section{WP2 — Synthesis of the Audiometric Spatialisation Test}
\label{sec:wp2}

\subsection{Objective}

The objective of WP2 was to identify, select, and implement a scientifically grounded
protocol for measuring a user's spatial hearing capabilities. The resulting test needed
to produce meaningful, interpretable key performance indicators (KPIs) while remaining
practical for participants interacting through a gamified web interface.

\subsection{Bibliographic Research (Task 2.1)}

A systematic review of the psychoacoustics literature was conducted to survey existing
methods for the objective assessment of sound localisation. Three broad families of
evaluation methods were identified.

\subsubsection{Interaural Time Difference and Interaural Level Difference Threshold Measurements}

ITD and ILD thresholds are classical measures of binaural hearing sensitivity.
ITD reflects the difference in arrival time between the two ears; ILD reflects the
difference in sound pressure level. Threshold estimation via adaptive procedures —
most commonly the \textit{staircase method} — converges on the smallest detectable
interaural difference for a given listener.

Although theoretically rigorous, the staircase paradigm was ultimately
\textbf{excluded} from consideration for this project. A standard two-down one-up
staircase with a termination criterion of ten reversals requires substantial trial
counts (typically 60–120 trials per condition), resulting in session durations that
are incompatible with the casual, engaging character that the gamified web application
is designed to maintain. Participant fatigue and attrition risk would undermine both
data quality and user retention. Furthermore, presenting sub-millisecond ITD stimuli
with the timing precision required by the staircase method imposes strict constraints
on audio hardware latency that cannot be reliably guaranteed across heterogeneous
consumer devices.

\subsubsection{Free-Field Localisation Tasks}

Free-field paradigms require participants to identify or reproduce the perceived
direction of a sound source. While ecologically valid, these tasks demand either
anechoic chamber conditions or carefully calibrated binaural rendering via
head-related transfer functions (HRTFs). Reliable HRTF-based spatialisation over
arbitrary consumer headphones introduces personalisation challenges — individual
anatomical differences cause significant listener-specific variation — making
standardised comparisons across participants difficult.

\subsubsection{The VASI Test}

The \textit{Virtual Auditory Space Identification} (VASI) test was selected as the
reference protocol for this project~\cite{vasi}. The VASI paradigm presents
participants with spatialised sound stimuli rendered through headphones using
binaural processing. Participants indicate the perceived location of each stimulus
from a set of discrete reference positions. The resulting confusion matrix and
derived metrics — including polar angle error, front–back confusion rate, and
lateral angle accuracy — provide a compact, multi-dimensional profile of spatial
hearing performance.

The VASI framework was preferred for the following reasons:

\begin{itemize}
    \item \textbf{Brevity}: A standard VASI session can be completed in 10–15 minutes,
    preserving participant engagement within a gamified context.
    \item \textbf{Richness of KPIs}: The protocol yields simultaneous metrics across
    elevation, azimuth, and front–back discrimination, capturing the multidimensional
    nature of spatial hearing in a single session.
    \item \textbf{Headphone compatibility}: Binaural rendering via generic HRTFs is
    sufficient for group-level comparisons and individual profiling, without requiring
    individualised measurement.
    \item \textbf{Prior validation}: VASI-type paradigms have been validated in
    peer-reviewed studies on both normal-hearing adults and populations with spatial
    hearing deficits, providing a normative reference for interpreting individual scores.
\end{itemize}

\subsection{Selection of Test Protocol and KPIs (Task 2.2)}

Based on the bibliographic review, the VASI protocol was adopted with the following
adaptations to suit the web deployment context:

\begin{itemize}
    \item Stimulus positions were drawn from a discrete set of virtual loudspeaker
    directions spanning the full sphere (front, rear, left, right, upper-front, upper-rear),
    rendered using open-source binaural impulse response datasets.
    \item Each trial presents a broadband noise burst (200~ms) spatialised to a randomly
    selected target direction. The participant selects the perceived direction from a
    graphical interface representing the spatial positions.
    \item The following KPIs are extracted from the response matrix:
    \begin{enumerate}
        \item \textbf{Lateral angle accuracy} (\%): percentage of trials in which the
        correct left–right hemisphere is identified.
        \item \textbf{Polar angle RMSE} (degrees): root mean square error between
        target and response elevation.
        \item \textbf{Front–back confusion rate} (\%): proportion of trials in which
        front and rear positions are exchanged.
        \item \textbf{Overall localisation score} (\%): composite metric weighting
        lateral accuracy, polar precision, and front–back discrimination.
    \end{enumerate}
\end{itemize}

\subsection{Implementation of the Test Script (Task 2.3)}

A standalone Python test script was developed to validate the stimulus generation pipeline
and scoring logic prior to web integration. The script is independent of WP1's
spatialisation engine (which handles real-time rendering); its role is to implement the
trial sequencing, response collection, and KPI computation logic.

Key implementation decisions include:

\begin{itemize}
    \item \textbf{Stimulus rendering}: The \texttt{soundfile} and \texttt{numpy} libraries
    are used to load, window, and convolve the noise bursts with the appropriate binaural
    room impulse responses (BRIRs).
    \item \textbf{Randomisation}: Trial order is fully randomised per session using a
    seeded pseudo-random number generator, ensuring reproducibility during validation.
    \item \textbf{Response interface}: In the standalone script, responses are collected
    via a command-line prompt; this interface is replaced by the web frontend in WP4.
    \item \textbf{KPI computation}: Confusion matrices are constructed using
    \texttt{scikit-learn}, and RMSE calculations are performed with \texttt{numpy}.
    Results are serialised to JSON for downstream reporting.
\end{itemize}

\subsection{Script Testing (Task 2.4)}

The Python script was validated through a combination of unit tests and pilot sessions.
Unit tests verified correctness of the confusion matrix construction, KPI formulas, and
edge cases (e.g., all-correct and all-incorrect response patterns). Pilot sessions with
three participants confirmed that session duration fell within the target range
(10–15 minutes) and that KPI distributions were plausible relative to published VASI
norms. Identified issues — primarily related to stimulus level calibration and trial
ordering artefacts — were corrected prior to integration with the web platform.

% ============================================================
%  WP4 — DEVELOPMENT OF THE WEB PLATFORM
% ============================================================
\section{WP4 — Development of the Web Platform}
\label{sec:wp4}

\subsection{Objective}

WP4 covers the full-stack design and implementation of the web platform that hosts the
spatialisation test developed in WP2. The platform makes the test accessible to end users
— competitive gamers and audio professionals — through a browser-based interface that
combines scientific rigour with an engaging, gamified user experience.

\subsection{Technology Stack and System Design (Task 4.1)}

\subsubsection{Architectural Overview}

The system follows a decoupled \textbf{client–server architecture} in which a single-page
application (SPA) frontend communicates with a RESTful backend API over HTTPS. This
separation of concerns offers several advantages: the frontend can be served as a static
bundle via a content delivery network (CDN), reducing latency for geographically distributed
users; the backend remains stateless and horizontally scalable; and the two layers can be
developed, tested, and deployed independently.

\begin{figure}[h]
    \centering
    % \includegraphics[width=0.85\textwidth]{figures/system_architecture.pdf}
    \caption{High-level system architecture of the spatialisation test platform.}
    \label{fig:architecture}
\end{figure}

\subsubsection{Frontend}

\textbf{React} (with TypeScript) was selected as the frontend framework on the following grounds:

\begin{itemize}
    \item \textbf{Component reusability}: The gamified test interface consists of discrete,
    reusable components (stimulus player, spatial response selector, progress indicator,
    results dashboard) that map naturally onto React's component model.
    \item \textbf{State management}: React's \texttt{useState}/\texttt{useReducer} hooks
    provide predictable, auditable state transitions across trial phases (idle → playing →
    awaiting response → feedback), which is essential for maintaining test integrity.
    \item \textbf{Web Audio API integration}: React's lifecycle model integrates cleanly
    with the \textit{Web Audio API}, which handles real-time binaural spatialisation of
    stimuli in the browser without server round-trips, thereby minimising latency.
    \item \textbf{Ecosystem maturity}: The availability of well-maintained libraries
    (e.g., \texttt{Howler.js} for audio scheduling, \texttt{Recharts} for results
    visualisation) reduced development overhead significantly.
\end{itemize}

\textbf{Tailwind CSS} was chosen for styling due to its utility-first approach, which
accelerates UI development and enforces visual consistency without requiring a bespoke
design system.

\subsubsection{Backend}

\textbf{FastAPI} (Python) was selected as the backend framework. The choice is motivated by:

\begin{itemize}
    \item \textbf{Language consistency}: The test logic developed in WP2 is written in
    Python; using FastAPI allows the KPI computation module to be imported directly into
    the API layer without reimplementation.
    \item \textbf{Automatic documentation}: FastAPI generates OpenAPI-compliant
    documentation from type annotations, facilitating integration testing and future
    third-party integrations.
    \item \textbf{Async support}: FastAPI's native \texttt{async}/\texttt{await} support
    ensures that I/O-bound operations (database writes, session lookups) do not block
    the event loop, supporting concurrent test sessions without additional infrastructure.
\end{itemize}

\textbf{PostgreSQL} was selected as the persistence layer for storing session metadata,
trial-level response data, and computed KPIs. Its ACID compliance guarantees that partial
session data cannot corrupt stored results, and its JSON column support allows flexible
storage of the trial response matrix without sacrificing query capability.

\subsubsection{Deployment}

The frontend is deployed as a static bundle; the backend runs in a containerised
environment (Docker). This configuration enables consistent behaviour across development
and production environments and simplifies horizontal scaling if user load increases.

\subsection{Backend API Layer (Task 4.2)}

The API exposes the following primary endpoints:

\begin{itemize}
    \item \texttt{POST /sessions} — Creates a new test session, returns a session token
    and the randomised trial sequence.
    \item \texttt{POST /sessions/\{id\}/responses} — Records the participant's response
    for a given trial; validates that the trial index is within the session scope and that
    the response is a valid spatial position identifier.
    \item \texttt{GET /sessions/\{id\}/results} — Triggers KPI computation (if not already
    cached) and returns the full results payload including the confusion matrix, per-KPI
    scores, and a percentile ranking derived from the normalised score distribution.
    \item \texttt{GET /leaderboard} — Returns an anonymised ranking of top scores,
    supporting the gamification layer.
\end{itemize}

Input validation is enforced via \textbf{Pydantic} models at every endpoint, ensuring
that malformed or out-of-range responses are rejected with informative error messages
before reaching the business logic layer.

\subsection{Frontend Development (Task 4.3)}

The frontend implements a three-phase user flow:

\begin{enumerate}
    \item \textbf{Onboarding}: The user is guided through a hardware calibration step
    (headphone check, volume normalisation via a reference tone) to ensure consistent
    stimulus presentation conditions.
    \item \textbf{Test execution}: Trials are presented sequentially. For each trial,
    the spatialised stimulus is played automatically via the Web Audio API; the user
    then selects the perceived direction from an interactive sphere visualisation.
    Real-time progress feedback (trial counter, animated transitions) sustains engagement.
    \item \textbf{Results presentation}: Upon completion, the user's KPI scores are
    displayed alongside a percentile ranking and radar chart comparing their performance
    across spatial dimensions (lateral, elevation, front–back). A leaderboard view
    provides social comparison, reinforcing the gamification element.
\end{enumerate}

Accessibility considerations include keyboard navigation for the spatial response
selector and ARIA labels on all interactive elements, ensuring the platform is usable
by participants with motor or visual impairments.

\subsection{Testing (Task 4.4)}

Testing was conducted at unit, integration, and end-to-end levels:

\begin{itemize}
    \item \textbf{Unit tests} (pytest, Vitest): Core functions — KPI computation,
    response validation, trial randomisation — were tested in isolation with
    parameterised inputs covering boundary conditions.
    \item \textbf{Integration tests}: API endpoint behaviour was verified against a
    test database instance, confirming correct session lifecycle management and
    result consistency across concurrent requests.
    \item \textbf{End-to-end tests} (Playwright): Critical user journeys (session
    creation → trial completion → results display) were automated across Chromium
    and Firefox to detect rendering and audio scheduling regressions.
    \item \textbf{Cross-device validation}: The platform was tested on representative
    desktop and mobile configurations to verify audio playback consistency and layout
    responsiveness.
\end{itemize}

All identified defects were resolved prior to the pilot deployment described in the
results section.
